\chapter{Strutture di comunicazione}
La comunicazione tra il client, il gestore dei messaggi e il server avvengono attraverso strutture json.
\section{Struttura JSON}
La struttura JSON che viene inviata per la comunicazione tra i vari elementi dovr\'a avere uno schema del tipo:
\begin{verbatim}
{
  msg_type: tipo_messaggio,
  msg_content: contenuto_del_messaggio
}
\end{verbatim}

I messaggi di ritorno avranno una struttura del tipo:

\begin{verbatim}
{
  result_type: tipo_messaggio_di_ritorno,
  result_content: contenuto_del_messaggio
}
\end{verbatim}

Codifichiamo quindi i messaggi per il gestore dei messaggi:

\begin{tabular}{|c|l|}
	\hline
	test\_messaggio & Serve per fare un semplice test di comunicazione\\
	\hline
	register\_station & Indica al sistema che la postazione vuole essere registrata sul sistema \\
	\hline
	unregister\_station & Indica al sistema che la postazione vuole essere tolta dal sistema \\
    \hline
	request\_code & Indica al sistema che la postazione vuole essere registrata sul sistema \\
    \hline
	register\_station & Indica al sistema che la postazione vuole essere registrata sul sistema \\
    \hline
\end{tabular}

